\chapter{Introduction}
\label{chap:intro}

\section{Framing}
\label{sec:amb}

This document describes the work developed in the context of the
final-year project of the Computer Science and Engineering programme of
the \ac{UBI}. The project was proposed by, and carried out at, the \ac{SEGAL}
laboratory of \ac{UBI}, under the supervision of Professor Paul Andrew
Crocker and the co-supervision of Vasco Senra, with day-to-day
infrastructure support from Luís Carvalho and Fernando Geraldes
(responsible for the laboratory's virtual machines, internal network
and \ac{VPN} access).

\ac{SEGAL} is the research laboratory of \ac{UBI} that operates and
maintains a \ac{GNSS} data and metadata service contributing to the
\ac{EPOS}, a pan-European research infrastructure focused on solid Earth
science~\cite{epos,eposGnss}. As part of the \ac{EPOS} Thematic Core
Service for \ac{GNSS} data and products, \ac{SEGAL} runs a Jakarta EE
\ac{REST} \ac{API} that exposes \ac{GNSS} station metadata and
observation files (in the \ac{RINEX} format~\cite{rinex}) to scientists,
operators and downstream services across Europe.

The \ac{API} in question has been in operation, in various forms, for
several years. Over that period the code base accumulated
technical debt: enourmous methods, and files of huge scale, duplicated logic copy-pasted across resources,
and a long tail of patches added by successive developers and
contributors to handle individual edge cases. A complete refactor of the
\ac{API} was therefore initiated, in which large parts of the
application were rewritten from scratch on top of Jakarta EE 11 and the
Payara Platform~\cite{jakartaee,payara}, while a small core of well
tested logic was carried over from the previous version.

A refactor of this magnitude poses a very specific risk: any seemingly
minor change in a query, a response field or a validation rule can
produce subtle, hard to detect regressions that only manifest when a
particular client, a particular station or a particular file is
requested. Until this project, validation of those changes was carried
out manually, a developer would deploy the new version, run a few
requests by hand, and compare the resulting payloads against the
previous one. This is what we are trying to solve.

\section{Motivation}
\label{sec:mot}

The motivation for this project is the gap between the pace at which
the new \ac{API} is being developed and the rate at which a human can
reasonably validate that each new version behaves exactly like the
previous one on the cases that already worked. Manual testing of a
\ac{REST} \ac{API} that returns thousands of records, in several
formats (\ac{JSON}, \ac{XML}, \ac{CSV}) and through dozens of query
combinations, is both error-prone and slow. This does not
scale to the rhythm of a continuous refactor where commits land
several times a day.

\medskip

The laboratory wanted an
\emph{automated} mechanism to compare every new build of the \ac{API}
against the previous one and flag behavioral divergences early, before
they reach the production deployment. On the other hand, the laboratory
also wanted to bring this mechanism inside a real \ac{CI}/\ac{CD}
pipeline running on infrastructure that the laboratory itself controls:
a free GitLab instance, paired with one or more GitLab
Runners hosted inside the laboratory's network. 

\medskip

The choice of hosting the
GitLab Runners was a necessity since the free tier only allows 400 compute
minutes, and due to the nature of app, and it's boot times, a full test suit 
usually takes up to 5 minutes per test, and being self-hosted allows us to 
connect to the real database infrastructure.
  

\medskip

Beyond the immediate value delivered to the \ac{EPOS} \ac{API}, designing and integrating this
framework was a significant personal milestone. It challenged me to collaborate with a team on
a live, rolling codebase—an experience that closely mirrors real-world software engineering.
Through this process, I gained practical, hands-on experience with modern Jakarta EE features,
\ac{CI}/\ac{CD} pipelines, self-hosted infrastructure, and regression testing against a moving
baseline, equipping me with industry-relevant skills that extend far beyond the scope of this
project.

\section{Objectives}
\label{sec:obj}

In line with the project proposal, the work pursues four objectives:
to design and implement an automated regression testing framework
for the Jakarta EE \ac{REST} \ac{API} on the Payara Platform, with
multi-format response comparison (\ac{JSON}, \ac{XML}, \ac{CSV}); to
integrate it into a self-managed GitLab \ac{CI}/\ac{CD} pipeline
running on \ac{VPN}-accessible runners; use of two prototype 
applications not only to validate the approach before applying it 
to the production API, but also to explore and learn about features 
that could be useful for the final implementation; and to produce 
structured per-commit comparisons at the level of \ac{HTTP}
responses, in a form consumable by both human developers and 
the \ac{CI} system.

The main contributions are a reusable three-stage GitLab pipeline
template for Jakarta EE applications, an \ac{HTTP}-level
differential test harness (\texttt{HttpDiff.java}) that boots one
Payara Micro per available \ac{WAR} and performs a three-way
comparison between \texttt{HEAD}, the immediate parent commit
\texttt{HEAD\textasciitilde 1} and a fixed reference commit
identified by the \texttt{DIFF\_FIXED\_BASELINE} \ac{CI}/\ac{CD}
variable, and two prototype Jakarta EE 11 applications retained as
didactic examples for future students.

\section{Organization of the Document}
\label{sec:organ}

The remaining chapters of this document are organized as follows.
\textbf{Chapter~\ref{chap:estado-da-arte}} reviews the state of the
art in test frameworks, \ac{CI}/\ac{CD} systems, and differential
testing. \textbf{Chapter~\ref{chap:tecno-ferra}} describes the
concrete technology stack adopted during the project.
\textbf{Chapter~\ref{chap:metodo-arq}} explains the iterative
methodology (two prototypes followed by the production application)
and the overall architecture of the testing framework.
\textbf{Chapter~\ref{chap:caso-epos}} applies the framework to the
production \ac{API} (\texttt{EPOS V4}) and walks through the
pipeline, the differential harness and the reports it produces.
\textbf{Chapter~\ref{chap:conc-trab-futuro}}
summarizes the results, reflects on what was learned and proposes
directions for future work.

\section{Repository Access}
\label{sec:repos}

The production \ac{API} (\texttt{GNSs V4}) is hosted in a private repository restricted to laboratory members \cite{repo-gnss-v4}. The two prototype applications developed during this work are publicly available in the author's personal namespace \cite{repo-java-ee, repo-syncspace}. 

Each repository contains a \texttt{README.md} with build instructions and (where applicable) the \texttt{.gitlab-ci.yml} configuration files discussed in Chapters~\ref{chap:metodo-arq} and~\ref{chap:caso-epos}.
